\begin{frame}{같은 EM, 다른 데이터: LDA 개요}
  GMM을 생성 과정 관점에서 보면, ``데이터가 \textbf{어디서}
  만들어졌는가''(잠재변수 $z_i$ 한 줄기)를 모델링. 그런데
  데이터가 \textbf{점}으로만 만들어지는 것은 아니다:
  \textbf{문서}(뉴스·논문·리뷰)의 경우 ``이 문서가 \textbf{어떤
  주제로} 만들어졌는가''를 물으면, 잠재변수는 \textbf{$K$개
  주제의 혼합 비율} $\theta_d$($\sum_k \theta_{dk} = 1$)가
  된다.
  \vskip0.6em
  이 질문을 \textbf{같은 EM의 틀}로 푸는 것이
  \kb{LDA}(Latent Dirichlet Allocation, 잠재 디리클레 할당)
  --- 2003년 \kb{블레이(Blei)·응(Ng)·조던(Jordan)}.
  \begin{itemize}\small
    \item 이름의 \textbf{Dirichlet} = 사전분포의 이름,
      \textbf{Allocation} = ``할당'' --- \textbf{문서가 토픽에
      할당되고, 토픽이 단어에 할당}된다는 생성 과정
    \item ``\textbf{GMM의 ``점 $\to$ 클러스터''를 ``단어
      $\to$ 토픽''으로 바꾼 것}'' --- 같은 EM 계열의 고전
      확률모델
  \end{itemize}
\end{frame}
